\documentclass[11pt]{article}
\usepackage[margin=1in]{geometry}
\usepackage{amsmath,amssymb,amsthm,hyperref}
\newtheorem{theorem}{Theorem}
\newcommand{\tr}{\operatorname{tr}}
\newcommand{\adj}{\operatorname{adj}}
\newcommand{\G}{\mathcal G}
\title{A quadratic counterexample to the product Hermitian Jacobi\linebreak eigenfunction conjecture\\\large MathDB 376316}
\author{Principia Math In Progress}
\date{17 September 2026}
\begin{document}
\maketitle

\paragraph{Task 1: the original problem.}
Conjecture 3.12 of \cite{K} asks whether every product polynomial defined
by equation (21) is an eigenfunction of the Hermitian Jacobi operator
on the matrix simplex (Definition 3.5). We give an interior counterexample
for two $2\times2$ matrices. The result is independent of the nonzero
normalization constants used for the Jacobi polynomials.

\begin{theorem}\label{thm}
Let $r,s>0$, $m=k=2$, and $\kappa=(r,s,s)$. The product polynomial with
$\tau_1=(0,0)$ and $\tau_2=(1,0)$ is not an eigenfunction of $\G^2_\kappa$.
In particular, Conjecture 3.12 is false for $\kappa=(1,1,1)$, which also
belongs to the geometric parameter range $n_1=n_2=n_3=m=2$.
\end{theorem}

\section*{Identifying the prescribed polynomial}
Write $A=\Lambda_1$, $B=\Lambda_2$, and $S=I_2-A$. The interior of the
matrix simplex is given by $A>0$, $B>0$, and $I_2-A-B>0$.
In equation (21) of \cite{K}, the chosen partitions give
$a_1=2+2s$ and $a_2=s$. The first Jacobi factor has partition zero and
is a nonzero constant. The second factor is, up to a nonzero constant,
\[
 P_{(1,0)}^{(s,s,2)}(Y)=\tr Y-1.
\]
Indeed, this is an invariant polynomial of degree one, orthogonal to
constants for the positive beta weight
$\det(Y)^{s-1}\det(I_2-Y)^{s-1}$ on $0<Y<I_2$. This weight is invariant
under $Y\mapsto I_2-Y$, so its mean trace is $1$. This uniquely determines
the affine polynomial up to a nonzero scalar. For $s=1$ the determinant
in Definition 1.9 gives the same conclusion directly: the divided
difference of $6u^2-6u+1$ at the two eigenvalues is
$6(\tr Y-1)$.

It follows that the specified product polynomial is a nonzero constant
multiple of
\begin{align}
 F(A,B)&=\det(S)\bigl(\tr(S^{-1/2}BS^{-1/2})-1\bigr)\notag\\
       &=\tr(\adj(S)B)-\det(S)\label{eq:F}\\
       &=\tr A+\tr B-1-\det A-M(A,B),\notag\\
 M(A,B)&=(\tr A)(\tr B)-\tr(AB).\notag
\end{align}
The last equality uses the elementary $2\times2$ identities
$\adj(I_2-A)=(1-\tr A)I_2+A$ and
$\det(I_2-A)=1-\tr A+\det A$. In particular, $F$ is a polynomial
in all matrix entries, invariant under simultaneous unitary conjugation.

\section*{Applying the full matrix operator}
Set $N=r+2s+3$, $X_1=A$, and $X_2=B$. Definition 3.5 specializes to
\begin{align}
 \G f={}&\sum_{j=1}^2\sum_{a,b,c,d=1}^2
 (I_2-X_j)_{ad}(X_j)_{cb}
 \frac{\partial^2f}{\partial(X_j)_{ab}\partial(X_j)_{cd}}\notag\\
 &-\frac12\sum_{j\ne\ell}\sum_{a,b,c,d=1}^2
 \bigl((X_j)_{ad}(X_\ell)_{cb}+(X_j)_{cb}(X_\ell)_{ad}\bigr)
 \frac{\partial^2f}{\partial(X_j)_{ab}\partial(X_\ell)_{cd}}\label{eq:G}\\
 &+\sum_{j=1}^2\sum_{a,b=1}^2
 \bigl((\kappa_j+1)\delta_{ab}-N(X_j)_{ab}\bigr)
 \frac{\partial f}{\partial(X_j)_{ab}}.\notag
\end{align}
Derivatives here are taken before restricting to any special matrices.
For example, when $A_{12}=u+iv$ and $A_{21}=u-iv$, they mean
$\partial_{A_{12}}=(\partial_u-i\partial_v)/2$ and
$\partial_{A_{21}}=(\partial_u+i\partial_v)/2$ on the Hermitian space.
In particular, setting off-diagonal entries to zero before differentiation
would give an incorrect operator calculation.

Direct differentiation in \eqref{eq:G} gives
\begin{align}
 \G\tr A&=2(r+1)-N\tr A,&
 \G\tr B&=2(s+1)-N\tr B,\label{eq:table}\\
 \G\det A&=r\tr A-2(N-1)\det A,\notag\\
 \G M&=(s+1)\tr A+(r+1)\tr B-2(N-1)M.\notag
\end{align}
Here the determinant identity follows by differentiating
$A_{11}A_{22}-A_{12}A_{21}$. For the last line, the two intermediate
identities are
\begin{align*}
 \G\bigl((\tr A)(\tr B)\bigr)
 &=2(r+1)\tr B+2(s+1)\tr A
   -2N(\tr A)(\tr B)-2\tr(AB),\\
 \G\tr(AB)
 &=(r+1)\tr B+(s+1)\tr A
   -2N\tr(AB)-2(\tr A)(\tr B).
\end{align*}
Subtracting proves the formula for $M$. Substitution into \eqref{eq:F}
now yields the polynomial identity
\begin{equation}\label{eq:residual}
 \boxed{\ \G F=-2(N-1)F+sL,\qquad
 L(A,B)=\tr A+2\tr B-2.\ }
\end{equation}

\section*{An exact interior witness and the spectral correction}
Take
\[
 A_0=\begin{pmatrix}1/3&0\\0&2/3\end{pmatrix},\qquad
 B_0=\frac29 I_2,\qquad
 I_2-A_0-B_0=\begin{pmatrix}4/9&0\\0&1/9\end{pmatrix}.
\]
All three matrices are positive definite. Also
$\det(I_2-A_0)=2/9$ and
$\tr(\adj(I_2-A_0)B_0)=2/9$, whence
\[
 F(A_0,B_0)=0,\qquad L(A_0,B_0)=-\frac19,\qquad
 (\G F)(A_0,B_0)=-\frac{s}{9}\ne0.
\]
If $F$ were an eigenfunction, $\G F=\lambda F$ for some constant
$\lambda$, the left side would vanish at every zero of $F$.
This contradiction proves Theorem \ref{thm}, including the stated
case $r=s=1$ with exact operator value $-1/9$.

The obstruction can also be described explicitly. Equations
\eqref{eq:table} imply $\G L=-NL$. Since $N>3$, the polynomial
\[
 Q=F-\frac{s}{N-2}L
 \quad\hbox{satisfies}\quad \G Q=-2(N-1)Q.
\]
Thus the proposed product $F$ combines two nonzero eigenfunctions,
$Q$ and $L$, with distinct eigenvalues $-2(N-1)$ and $-N$.
An overall rescaling of the operator or of the Jacobi polynomials
does not remove the counterexample. No boundary parameter, singular
matrix, or omitted regularity assumption is involved.

\paragraph{Verification scope.}
The preceding identities are the proof. The accompanying exact checker
independently differentiates the source operator in all eight formal
matrix entries and then in eight real Hermitian coordinates. It also
checks the rational witness, the affine Jacobi normalization, and the
two-eigenfunction decomposition. No floating-point inference is used.
The proof and source audit were performed by the authoring agent;
independent peer review is not claimed.

\section*{Four follow-up tasks}
\paragraph{Task 2.}
Classify which product indices of total polynomial degree at most four
are genuine eigenfunctions for $m=k=2$, including exceptional parameter
sets where distinct eigenvalues coincide.
\paragraph{Task 3.}
Develop a finite triangular procedure that decomposes each product
polynomial into true Jacobi eigenfunctions, and determine whether its
coefficients admit formulas in terms of partitions and the parameters.
\paragraph{Task 4.}
Resolve the separate strong-triangularity Conjecture 3.13, keeping track
of multiplicity spaces of the multivariate Schur basis. The present
counterexample does not decide that conjecture.
\paragraph{Task 5.}
For general matrix size $m$ and $k\ge2$, determine the first product
indices that fail the eigenfunction property and construct explicit
interior witnesses or prove any surviving families.

\begin{thebibliography}{9}
\bibitem{K} T. Kuijper, \emph{Jacobi polynomials on the simplex of Hermitian
matrices and zonal spherical functions on partial flag manifolds},
\href{https://arxiv.org/abs/2607.25760v1}{arXiv:2607.25760v1} (2026),
Definition 1.9, equation (7), Definition 3.5, equation (21), and
Conjecture 3.12 (pp.~10--11, 30, 34--35).
\end{thebibliography}
\end{document}
